\documentclass[11pt,a4paper]{article}
\usepackage{fontspec,amsmath,amssymb,amsthm,microtype}
\setmainfont{lmroman10-regular.otf}[BoldFont=lmroman10-bold.otf,ItalicFont=lmroman10-italic.otf,BoldItalicFont=lmroman10-bolditalic.otf]
\usepackage{xcolor}
\usepackage[most]{tcolorbox}
\usepackage{mathrsfs}
\providecommand{\lt}{<}
\providecommand{\gt}{>}
\setmonofont{DejaVuSansMono.ttf}[BoldFont=DejaVuSansMono-Bold.ttf,ItalicFont=DejaVuSansMono-Oblique.ttf,BoldItalicFont=DejaVuSansMono-BoldOblique.ttf,Scale=0.86]
\newfontfamily\leanfallback{STIXTwoMath-Regular.otf}[Scale=0.86]
\definecolor{leanrule}{gray}{0.62}
\definecolor{leanback}{gray}{0.962}
\definecolor{leantitle}{gray}{0.22}
\definecolor{leanmuted}{gray}{0.38}
\newtcolorbox{leanbox}[1]{enhanced,breakable,colback=leanback,colframe=leanrule,
  boxrule=0.4pt,arc=0.6mm,left=2.4mm,right=2.4mm,top=1.8mm,bottom=1.8mm,
  before skip=0.8em,after skip=1em,fonttitle=\small\bfseries,coltitle=leantitle,
  colbacktitle=leanback,titlerule=0.3pt,title={#1}}
\newlength{\leancharwidth}
\newenvironment{leancode}{\par\smallskip\begingroup\ttfamily\footnotesize\raggedright
  \settowidth{\leancharwidth}{M}\setlength{\parskip}{0pt}\setlength{\parindent}{0pt}
  \frenchspacing}{\par\endgroup\smallskip}
\newcommand{\leanline}[2]{\hangindent=\dimexpr#1\leancharwidth+4\leancharwidth\relax
  \hangafter=1\noindent\hspace*{#1\leancharwidth}#2\par}
\newcommand{\leansource}[1]{\par{\footnotesize\color{leanmuted}#1}\par}
\newcommand{\leanlabel}[1]{\par\smallskip\noindent{\small\bfseries #1}\par}
\usepackage[margin=25mm]{geometry}
\usepackage[unicode,hidelinks]{hyperref}
\hypersetup{pdfauthor={LangevinLean Project}}
\setlength{\emergencystretch}{3em}
\setlength{\parskip}{0.4em}
\setlength{\parindent}{0pt}
\allowdisplaybreaks[2]
\raggedbottom
\title{\texorpdfstring{Local error relative to exact Hamiltonian motion}{Local error relative to exact Hamiltonian motion}}
\author{LangevinLean Project}
\date{}
\begin{document}\maketitle
Draft mathematical exposition; these notes may contain mistakes.


To quantify sampling bias, compare the defined \(O\)–Verlet–\(O\) step with a reference that uses exact Hamiltonian motion between the same two thermostat updates. Exact Hamiltonian motion and thermostat preserve the Gibbs target separately. This section derives the local state error from the Hamiltonian integral equations. It works in any real Hilbert space, including \(L^{2}\), where uniform velocity and force norm bounds are second-moment statements.



The defect is derived by integrating the true Hamiltonian equations. A common-noise comparison then inserts this deterministic estimate into the same weighted geometry used for OBABO contraction.


\section{\texorpdfstring{Integrate the Hamiltonian equations instead of assuming a local error}{Integrate the Hamiltonian equations instead of assuming a local error}}


The Hamiltonian equations are \(v(t)-v(0)=-\int_0^t g(x(s))ds\) and \(x(t)-x(0)=\int_0^t v(s)ds\). If the force norm is at most \(G\), the integral norm inequality gives \(\|v(t)-v(0)\|\le Gt\). Subtract \(tv(0)\) from the position equation; the result is the integral of the velocity increment. Integrating the derived velocity bound \(Gs\) gives the Euler-position defect \(\|x(t)-x(0)-tv(0)\|\le Gt^2/2\). One proves the exact integral identities using integrability from continuity before applying either inequality.



The exact Hamiltonian solution satisfies integrated velocity and position equations. Substitute the position primitive into the force difference and use smoothness to obtain explicit integral remainders. This derives a local state error from the differential equation itself.

\begin{leanbox}{Lean: \texttt{hamiltonian\_\allowbreak{}velocity\_\allowbreak{}increment\_\allowbreak{}bound}, \texttt{hamiltonian\_\allowbreak{}position\_\allowbreak{}euler\_\allowbreak{}defect}}
\leansource{\texttt{hamiltonian\_\allowbreak{}velocity\_\allowbreak{}increment\_\allowbreak{}bound} \quad KineticLangevin/\allowbreak{}VerletLocal.lean, lines 20–29}
\begin{leancode}
\leanline{0}{theorem\ hamiltonian\_velocity\_increment\_bound\ (g\ :\ E\ →\ E)\ (v\ x\ :\ ℝ\ →\ E)\ (T\ G\ :\ ℝ)}
\leanline{4}{(hbound\ :\ ∀\ s\ ∈\ Icc\ 0\ T,\ ‖g\ (x\ s)‖\ ≤\ G)}
\leanline{4}{(heq\ :\ ∀\ t\ ∈\ Icc\ 0\ T,\ v\ t\ -\ v\ 0\ =\ -(∫\ s\ in\ 0..t,\ g\ (x\ s)))}
\leanline{4}{(t\ :\ ℝ)\ (ht\ :\ t\ ∈\ Icc\ 0\ T)\ :\ ‖v\ t\ -\ v\ 0‖\ ≤\ G\ *\ t}
\end{leancode}
\leansource{\texttt{hamiltonian\_\allowbreak{}position\_\allowbreak{}euler\_\allowbreak{}defect} \quad KineticLangevin/\allowbreak{}VerletLocal.lean, lines 31–49}
\begin{leancode}
\leanline{0}{theorem\ hamiltonian\_position\_euler\_defect\ (g\ :\ E\ →\ E)\ (v\ x\ :\ ℝ\ →\ E)\ (T\ G\ :\ ℝ)}
\leanline{4}{(hv\ :\ Continuous\ v)}
\leanline{4}{(hbound\ :\ ∀\ s\ ∈\ Icc\ 0\ T,\ ‖g\ (x\ s)‖\ ≤\ G)}
\leanline{4}{(heq\ :\ ∀\ t\ ∈\ Icc\ 0\ T,\ v\ t\ -\ v\ 0\ =\ -(∫\ s\ in\ 0..t,\ g\ (x\ s)))}
\leanline{4}{(hkin\ :\ ∀\ t\ ∈\ Icc\ 0\ T,\ x\ t\ -\ x\ 0\ =\ ∫\ s\ in\ 0..t,\ v\ s)}
\leanline{4}{(t\ :\ ℝ)\ (ht\ :\ t\ ∈\ Icc\ 0\ T)\ :\ ‖x\ t\ -\ (x\ 0\ +\ t\ •\ v\ 0)‖\ ≤\ G\ *\ t\ \textasciicircum{}\ 2\ /\ 2}
\end{leancode}
\leanlabel{How the Lean proof follows the argument}
\texttt{hamiltonian\_\allowbreak{}velocity\_\allowbreak{}increment\_\allowbreak{}bound} gives the formal counterpart of this step. The main intermediate facts are \texttt{hb} in that order. Named proof applications include \texttt{hbound}; the adjacent proof shows their hypotheses and order of use.




\texttt{hamiltonian\_\allowbreak{}position\_\allowbreak{}euler\_\allowbreak{}defect} gives the formal counterpart of this step. The main intermediate facts are \texttt{hid} in that order. Named proof applications include \texttt{interval\allowbreak{}Integral.\allowbreak{}norm\_\allowbreak{}integral\_\allowbreak{}le\_\allowbreak{}of\_\allowbreak{}norm\_\allowbreak{}le}; the adjacent proof shows their hypotheses and order of use.
\end{leanbox}

\section{\texorpdfstring{Compare both Verlet components to the true endpoint}{Compare both Verlet components to the true endpoint}}


The defined drift-stage position is \(x_V=x_0+hv_0-h^2g(x_0)/2\). Its difference from the Euler position has norm at most \(Gh^2/2\); adding the exact position defect proves \(\|x_V-x(h)\|\le Gh^2\). For velocity, subtract the frozen-force Euler value \(v_0-hg(x_0)\) from both trajectories. The exact difference is the integral of \(g(x(s))-g(x_0)\), bounded by \(MC h^2/2\) using Lipschitzness and the kinematics with velocity norm \(C\). The Verlet difference is exactly \(-h(g(x_V)-g(x_0))/2\), whose norm is at most \(Mh^2C/2+MGh^3/4\). Combining the two gives \(\|v_V-v(h)\|\le MC h^2+MGh^3/4\). Every residual is computed from the update and the exact path; none is supplied as an additional accuracy assumption.



The first half kick, drift, and final half kick are compared componentwise with those exact integrals. Terms sharing the same leading Taylor contribution cancel. The remaining bounds retain the initial velocity and force magnitudes needed for subsequent moment integration.

\begin{leanbox}{Lean: \texttt{verlet\_\allowbreak{}local\_\allowbreak{}component\_\allowbreak{}bounds}}
\leansource{\texttt{verlet\_\allowbreak{}local\_\allowbreak{}component\_\allowbreak{}bounds} \quad KineticLangevin/\allowbreak{}VerletLocal.lean, lines 51–110}
\begin{leancode}
\leanline{0}{theorem\ verlet\_local\_component\_bounds\ (g\ :\ E\ →\ E)\ (M\ :\ ℝ≥0)\ (hg\ :\ LipschitzWith\ M\ g)}
\leanline{4}{(v\ x\ :\ ℝ\ →\ E)\ (h\ C\ G\ :\ ℝ)\ (hh\ :\ 0\ ≤\ h)\ (hC\ :\ 0\ ≤\ C)\ (hG\ :\ 0\ ≤\ G)}
\leanline{4}{(hv\ :\ Continuous\ v)\ (hx\ :\ Continuous\ x)}
\leanline{4}{(hvelocity\ :\ ∀\ s\ ∈\ Icc\ 0\ h,\ ‖v\ s‖\ ≤\ C)}
\leanline{4}{(hforce\ :\ ∀\ s\ ∈\ Icc\ 0\ h,\ ‖g\ (x\ s)‖\ ≤\ G)}
\leanline{4}{(heq\ :\ ∀\ t\ ∈\ Icc\ 0\ h,\ v\ t\ -\ v\ 0\ =\ -(∫\ s\ in\ 0..t,\ g\ (x\ s)))}
\leanline{4}{(hkin\ :\ ∀\ t\ ∈\ Icc\ 0\ h,\ x\ t\ -\ x\ 0\ =\ ∫\ s\ in\ 0..t,\ v\ s)\ :}
\leanline{4}{‖(verletStep\ g\ h\ (v\ 0,x\ 0)).2\ -\ x\ h‖\ ≤\ G\ *\ h\ \textasciicircum{}\ 2\ ∧}
\leanline{6}{‖(verletStep\ g\ h\ (v\ 0,x\ 0)).1\ -\ v\ h‖\ ≤\ (M\ :\ ℝ)\ *\ C\ *\ h\ \textasciicircum{}\ 2\ +\ (M\ :\ ℝ)\ *\ G\ *\ h\ \textasciicircum{}\ 3\ /\ 4}
\end{leancode}
\leanlabel{How the Lean proof follows the argument}
\texttt{verlet\_\allowbreak{}local\_\allowbreak{}component\_\allowbreak{}bounds} gives the formal counterpart of this step. The main intermediate facts are \texttt{h0}, \texttt{hhmem}, \texttt{h\allowbreak{}V0}, \texttt{h\allowbreak{}F0}, \texttt{xe}, \texttt{xn} in that order. Named proof applications include \texttt{mul\_\allowbreak{}le\_\allowbreak{}mul\_\allowbreak{}of\_\allowbreak{}nonneg\_\allowbreak{}left}, \texttt{velocity\allowbreak{}Defect\_\allowbreak{}bound}; the adjacent proof shows their hypotheses and order of use.
\end{leanbox}

\section{\texorpdfstring{Express the state defect in the contracting OBABO metric}{Express the state defect in the contracting OBABO metric}}


The upper metric equivalence and completing-square Hilbert representation imply the convenient component estimate \(\|(v,x)\|_E\le2(\|x\|+\|v\|/\sqrt M)\). The proof squares both nonnegative sides, expands the scalar square, and uses \(\sqrt M^2=M\). Inserting the two component defects yields

\[\|\mathrm{Verlet}_h(v_0,x_0)-(v(h),x(h))\|_E\le2h^2\bigl(G+\sqrt M C+\sqrt M G h/4\bigr).\]

This estimate uses only a Lipschitz force and the explicit moment bounds; no Hessian regularity is required.



The physical component errors are inserted into the verified weighted transformation. The triangle inequality yields a phase defect in the same norm used for contraction. Keeping this common metric avoids an extra norm-conversion loss at each later iteration.

\begin{leanbox}{Lean: \texttt{weightedPhaseNorm\_\allowbreak{}le\_\allowbreak{}components}, \texttt{verlet\_\allowbreak{}local\_\allowbreak{}weighted\_\allowbreak{}bound}}
\leansource{\texttt{weightedPhaseNorm\_\allowbreak{}le\_\allowbreak{}components} \quad KineticLangevin/\allowbreak{}VerletLocal.lean, lines 112–128}
\begin{leancode}
\leanline{0}{theorem\ weightedPhaseNorm\_le\_components\ (γ\ L\ :\ ℝ)\ (hγ\ :\ 0\ <\ γ)\ (hL\ :\ 0\ <\ L)}
\leanline{4}{(hfriction\ :\ 4\ *\ L\ ≤\ γ\ \textasciicircum{}\ 2)\ (s\ :\ E\ ×\ E)\ :}
\leanline{4}{weightedPhaseNorm\ (1\ /\ L)\ (1\ /\ γ)\ s\ ≤\ 2\ *\ (‖s.2‖\ +\ ‖s.1‖\ /\ Real.sqrt\ L)}
\end{leancode}
\leansource{\texttt{verlet\_\allowbreak{}local\_\allowbreak{}weighted\_\allowbreak{}bound} \quad KineticLangevin/\allowbreak{}VerletLocal.lean, lines 130–155}
\begin{leancode}
\leanline{0}{theorem\ verlet\_local\_weighted\_bound\ (g\ :\ E\ →\ E)\ (M\ :\ ℝ≥0)\ (hg\ :\ LipschitzWith\ M\ g)}
\leanline{4}{(v\ x\ :\ ℝ\ →\ E)\ (γ\ h\ C\ G\ :\ ℝ)\ (hγ\ :\ 0\ <\ γ)\ (hM\ :\ 0\ <\ (M\ :\ ℝ))}
\leanline{4}{(hfriction\ :\ 4\ *\ (M\ :\ ℝ)\ ≤\ γ\ \textasciicircum{}\ 2)\ (hh\ :\ 0\ ≤\ h)\ (hC\ :\ 0\ ≤\ C)\ (hG\ :\ 0\ ≤\ G)}
\leanline{4}{(hv\ :\ Continuous\ v)\ (hx\ :\ Continuous\ x)}
\leanline{4}{(hvelocity\ :\ ∀\ s\ ∈\ Icc\ 0\ h,\ ‖v\ s‖\ ≤\ C)}
\leanline{4}{(hforce\ :\ ∀\ s\ ∈\ Icc\ 0\ h,\ ‖g\ (x\ s)‖\ ≤\ G)}
\leanline{4}{(heq\ :\ ∀\ t\ ∈\ Icc\ 0\ h,\ v\ t\ -\ v\ 0\ =\ -(∫\ s\ in\ 0..t,\ g\ (x\ s)))}
\leanline{4}{(hkin\ :\ ∀\ t\ ∈\ Icc\ 0\ h,\ x\ t\ -\ x\ 0\ =\ ∫\ s\ in\ 0..t,\ v\ s)\ :}
\leanline{4}{weightedPhaseNorm\ (1\ /\ (M\ :\ ℝ))\ (1\ /\ γ)\ (verletStep\ g\ h\ (v\ 0,x\ 0)\ -\ (v\ h,x\ h))\ ≤}
\leanline{6}{2\ *\ h\ \textasciicircum{}\ 2\ *\ (G\ +\ Real.sqrt\ (M\ :\ ℝ)\ *\ C\ +\ Real.sqrt\ (M\ :\ ℝ)\ *\ G\ *\ h\ /\ 4)}
\end{leancode}
\leanlabel{How the Lean proof follows the argument}
\texttt{weighted\allowbreak{}Phase\allowbreak{}Norm\_\allowbreak{}le\_\allowbreak{}components} gives the formal counterpart of this step. The main intermediate facts are \texttt{hmetric}, \texttt{hab}, \texttt{hu}, \texttt{hn}, \texttt{hsq}, \texttt{hp} in that order. Its remaining equalities are discharged by the simplifications displayed in the proof.




\texttt{verlet\_\allowbreak{}local\_\allowbreak{}weighted\_\allowbreak{}bound} gives the formal counterpart of this step. The main intermediate facts are \texttt{he}, \texttt{hroot}, \texttt{hratio} in that order. The final step is the algebraic normalization or ordered arithmetic displayed in the proof.
\end{leanbox}

\section{\texorpdfstring{Couple the defined OBABO step with the exact target-preserving split}{Couple the defined OBABO step with the exact target-preserving split}}


Start the Hamiltonian path at the first Ornstein–Uhlenbeck half-step. Apply the same final thermostat vector to its endpoint and to the Verlet endpoint. The final additive noise cancels exactly, while the velocity defect is multiplied by \(\delta\). Since \(|\delta|\le1\), its norm is no larger than the Verlet velocity defect, and the position defect is unchanged. Applying the component metric bound therefore gives the same explicit local estimate for the \(O\)–Verlet–\(O\) and \(O\)–exact-Hamiltonian–\(O\) outputs. In \(L^{2}\), an invariant Gibbs Hamiltonian reference supplies \(C=\sqrt d\) and \(G\le\sqrt{Md}\) from the Gaussian velocity marginal and Gibbs score moment theorem. The local bound itself keeps the exact path equations and norm bounds explicit so the probabilistic reference construction must verify them.



The defined OBABO step and its exact comparator use identical first and second thermostat noises. Their only difference is replacing the central exact Hamiltonian evolution by Verlet. Thus the deterministic Hamiltonian local defect directly controls the coupled stochastic step discrepancy.

\begin{leanbox}{Lean: \texttt{obabo\_\allowbreak{}hamiltonian\_\allowbreak{}local\_\allowbreak{}bound}}
\leansource{\texttt{obabo\_\allowbreak{}hamiltonian\_\allowbreak{}local\_\allowbreak{}bound} \quad KineticLangevin/\allowbreak{}VerletLocal.lean, lines 157–195}
\begin{leancode}
\leanline{0}{theorem\ obabo\_hamiltonian\_local\_bound\ (g\ :\ E\ →\ E)\ (M\ :\ ℝ≥0)\ (hg\ :\ LipschitzWith\ M\ g)}
\leanline{4}{(v\ x\ :\ ℝ\ →\ E)\ (γ\ h\ δ\ σ\ C\ G\ :\ ℝ)\ (hγ\ :\ 0\ <\ γ)\ (hM\ :\ 0\ <\ (M\ :\ ℝ))}
\leanline{4}{(hfriction\ :\ 4\ *\ (M\ :\ ℝ)\ ≤\ γ\ \textasciicircum{}\ 2)\ (hh\ :\ 0\ ≤\ h)\ (hC\ :\ 0\ ≤\ C)\ (hG\ :\ 0\ ≤\ G)}
\leanline{4}{(hδ\ :\ |δ|\ ≤\ 1)\ (hv\ :\ Continuous\ v)\ (hx\ :\ Continuous\ x)}
\leanline{4}{(hvelocity\ :\ ∀\ s\ ∈\ Icc\ 0\ h,\ ‖v\ s‖\ ≤\ C)}
\leanline{4}{(hforce\ :\ ∀\ s\ ∈\ Icc\ 0\ h,\ ‖g\ (x\ s)‖\ ≤\ G)}
\leanline{4}{(heq\ :\ ∀\ t\ ∈\ Icc\ 0\ h,\ v\ t\ -\ v\ 0\ =\ -(∫\ s\ in\ 0..t,\ g\ (x\ s)))}
\leanline{4}{(hkin\ :\ ∀\ t\ ∈\ Icc\ 0\ h,\ x\ t\ -\ x\ 0\ =\ ∫\ s\ in\ 0..t,\ v\ s)}
\leanline{4}{(z\ :\ E\ ×\ E)\ (ξ\ ξ'\ :\ E)\ (hinit\ :\ (v\ 0,x\ 0)\ =\ ornsteinStep\ δ\ σ\ ξ\ z)\ :}
\leanline{4}{weightedPhaseNorm\ (1\ /\ (M\ :\ ℝ))\ (1\ /\ γ)}
\leanline{6}{(obaboStep\ g\ h\ δ\ σ\ z\ ξ\ ξ'\ -\ ornsteinStep\ δ\ σ\ ξ'\ (v\ h,x\ h))\ ≤}
\leanline{6}{2\ *\ h\ \textasciicircum{}\ 2\ *\ (G\ +\ Real.sqrt\ (M\ :\ ℝ)\ *\ C\ +\ Real.sqrt\ (M\ :\ ℝ)\ *\ G\ *\ h\ /\ 4)}
\end{leancode}
\leanlabel{How the Lean proof follows the argument}
\texttt{obabo\_\allowbreak{}hamiltonian\_\allowbreak{}local\_\allowbreak{}bound} gives the formal counterpart of this step. The main intermediate facts are \texttt{hc}, \texttt{hid}, \texttt{hp}, \texttt{hvbound}, \texttt{hroot}, \texttt{hratio} in that order. Named proof applications include \texttt{hc.\allowbreak{}1}; the adjacent proof shows their hypotheses and order of use.
\end{leanbox}

\end{document}
